\subsection{互补子模}

定义 8. 在 A-模 E 中，称两个子模 \(M_{1}\) , \(M_{2}\) 是互补的，如果 E 是 \(M_{1}\) 与 \(M_{2}\) 的直和。

第 8 号的命题 11 表明，要使 \(\mathbf{M}_1\) 与 \(\mathbf{M}_2\) 互补，必须且只需 \(\mathbf{M}_1 + \mathbf{M}_2 = \mathbf{E}\) 且 \(\mathbf{M}_1 \cap \mathbf{M}_2 = \{0\}\) （参见 I，§4，第 9 号，命题 15）。

命题 13. 设 \(\mathbf{M}_1, \mathbf{M}_2\) 是 A-模 E 中的两个互补子模。在 \(\mathbf{M}_1\) 上，典范映射 \(\mathbf{E} \to \mathbf{E} / \mathbf{M}_2\) 的限制是 \(\mathbf{M}_1\) 到 \(\mathbf{E} / \mathbf{M}_2\) 上的同构。

这个线性映射是满射，因为 \(M_{1} + M_{2} = E\) ；它也是单射，因为它的核是 \(M_{1}\) 与 \(M_{2}\) （ \(E \rightarrow E/M_{2}\) 的核）的交，从而约化为 \(\{0\}\) 。

推论。若 \(M_{2}\) 和 \(M_{2}^{\prime}\) 是 E 的同一子模 \(M_{1}\) 的两个补，则由有序对 \((x, x^{\prime}) \in \mathbf{M}_{2} \times \mathbf{M}_{2}^{\prime}\) （满足 \(x - x^{\prime} \in M_{1}\) ）组成的集合是 \(M_{2}\) 到 \(M_{2}^{\prime}\) 上的一个同构的图像。

立即可以看出，它是复合同构 \(M_{2} \rightarrow E/M_{1} \rightarrow M_{2}'\) 的图像。

定义 9. A-模 E 的一个子模 M 称为 E 的直因子，如果它在 E 中有一个互补子模。

此时，E/M 同构于 M 的一个补（命题 13）。

一个子模不一定有补（习题 11）。当一个子模是直因子时，它一般有若干个不同的补；但这些补彼此典范同构（命题 13 的推论）。

命题 14. 要使模 E 的子模 M 是直因子，必须且只需存在 E 的一个投影算子其像为 M，或存在 E 的一个投影算子其核为 M。

这由第 8 号的命题 12 及其推论立即得出。

命题 15. 给定 A-模的一个正合列

\[
0 \longrightarrow \mathrm{E} \stackrel {f} {\longrightarrow} \mathrm{F} \stackrel {g} {\longrightarrow} \mathrm{G} \longrightarrow 0\tag{33}
\]

下列命题等价：

\begin{enumerate}
\def\labelenumi{(\alph{enumi})}
\item
  子模 \(f(\mathbf{E})\) 在 \(\mathbf{F}\) 中是直因子。
\item
  存在线性收缩 \(r: \mathbf{F} \to \mathbf{E}\) 与 \(f\) 相伴（集合论，II，§3，第 8 号，定义 11）。
\item
  存在线性截面 \(s: \mathbf{G} \to \mathbf{F}\) 与 \(g\) 相伴（集合论，II，§3，第 8 号，定义 11）。
\end{enumerate}

此时， \(f + s: \mathbf{E} \oplus \mathbf{G} \to \mathbf{F}\) 是同构。

若存在投影算子 \(e\) 属于 \(\operatorname{End}(\mathbf{F})\) ，使得 \(e(\mathbf{F}) = f(\mathbf{E})\) ，则同态 \(f^{-1} \circ e: \mathbf{F} \to \mathbf{E}\) 是与 \(f\) 相伴的线性收缩；反之，若存在这样的收缩 \(r\) ，则立即可以看出 \(f \circ r\) 是 \(\mathbf{F}\) 中的一个投影算子，其像为 \(f(\mathbf{E})\) ，故 (a) 与 (b) 等价（命题 14）。若 \(f(\mathbf{E})\) 有一个补 \(\mathbf{E}'\) （在 \(\mathbf{F}\) 中），且 \(j: \mathbf{E}' \to \mathbf{F}\) 是典范单射，则 \(g \circ j\) 是 \(\mathbf{E}'\) 到 \(\mathbf{G}\) 上的同构，其逆同构（视为 \(\mathbf{G}\) 到 \(\mathbf{F}\) 中的映射）是与 \(g\) 相伴的线性截面。反之，若这样的截面 \(s\) 存在，则 \(s \circ g\) 是 \(\mathbf{F}\) 中的一个投影算子，其核为 \(f(\mathbf{E})\) ，故 (a) 与 (c) 等价（命题 14）。此外， \(s\) 是 \(\mathbf{G}\) 到 \(s(\mathbf{G})\) 上的双射，且由于 \(s(\mathbf{G})\) 与 \(f(\mathbf{E}), f + s\) 互补，后者是同构。

注意，给定 r（相应地，s）等价于给定 \(f(\mathrm{E})\) 在 F 中的一个补，即 r 的核（相应地，s 的像）。

当正合列 (33) 满足命题 15 的条件时，称它分裂，或称 \((\mathbf{F}, f, g)\) 是 \(\mathbf{G}\) 借助 \(\mathbf{E}\) 的平凡扩张（I，§6，第 1 号）。

推论 1. 设 \(u: \mathbf{E} \to \mathbf{F}\) 是线性映射。要存在线性映射 \(v: \mathbf{F} \to \mathbf{E}\) 使得 \(u \circ v = 1_{\mathbf{F}}\) （此时称 \(u\) 是右可逆的，而 \(v\) 称为 \(u\) 的右逆），必须且只需 \(u\) 是满射且其核是 \(\mathbf{E}\) 中的直因子。此时子模 \(\operatorname{Im}(v)\) （ \(\mathbf{E}\) 的子模）是 \(\operatorname{Ker}(u)\) 的一个补。

显然必须 \(u\) 是满射；而此时 \(v\) 是与 \(u\) 相伴的截面，结论由命题 15 得出。

推论 2. 设 \(u: \mathbf{E} \to \mathbf{F}\) 是线性映射。要存在线性映射 \(v: \mathbf{F} \to \mathbf{E}\) 使得 \(v \circ u = 1_{\mathbf{E}}\) （此时称 \(u\) 是左可逆的，而 \(v\) 称为 \(u\) 的左逆），必须且只需 \(u\) 是单射且其像是 \(\mathbf{F}\) 中的直因子。此时子模 \(\operatorname{Ker}(v)\) （ \(\mathbf{F}\) 的子模）是 \(\operatorname{Im}(u)\) 的一个补。

显然必须 u 是单射；而此时 v 是与 u 相伴的收缩，结论也由命题 15 得出。

注记。(1) 设 M, N 是 A-模 E 中的两个互补子模，p, q 分别是 E 到 M 上和到 N 上的、对应于 E 分解为 M 与 N 的直和的投影算子。已知（第 6 号，命题 6 的推论 1），对每个 A-模 F，映射 \((u, v) \mapsto u \circ p + v \circ q\) 是

\[
\operatorname{Hom} _ {\mathbf {A}} (\mathbf {M}, \mathbf {F}) \oplus \operatorname{Hom} _ {\mathbf {A}} (\mathbf {N}, \mathbf {F})
\]

到 \(\operatorname{Hom}_{\mathbf{A}}(\mathbf{E},\mathbf{F})\) 上的同构。 \(\operatorname{Hom}_{\mathbf{A}}(\mathbf{M},\mathbf{F})\) 在这个同构下的像是线性映射 \(w:\mathrm{E}\to \mathrm{F}\) （满足 \(w(x) = 0\) ，对所有 \(x\in N\) ）的集合。

\begin{enumerate}
\def\labelenumi{(\arabic{enumi})}
\setcounter{enumi}{1}
\tightlist
\item
  若 M, N 是 E 的两个子模，使得 \(M \cap N\) 是 M 与 N 的直因子，则 \(M \cap N\) 也是 \(M + N\) 的直因子：若 P（相应地，Q）是 \(M \cap N\) 在 M（相应地，N）中的一个补，则 \(M + N\) 是 \(M \cap N\) 、P 与 Q 的直和，这可以立即验证。
\end{enumerate}

